% Input fragment: exact three-coordinate charts for quadratic-field points.

\section{Three rational coordinates for an elliptic point over a quadratic field}
\label{sec:quadratic-point-compression}

The following elementary elimination reduces the ambient coordinate count
of a point over a quadratic field.  It trades degree for variables, which
is appropriate when the number of variables is the first optimization
criterion.  The square-root expressions below are reconstructed rational
functions, not new quantified variables.

\begin{proposition}[Eliminating one coordinate of a quadratic square root]
\label{prop:quadratic-square-root}
Let $F$ be a field of characteristic different from two, let
$\delta\in F^\times$ be a nonsquare, and write $s^2=\delta$.
Fix $j\in F$ and put $\lambda=j^2-\delta\ne0$.
For $A,B,C\in F$ with $C\ne0$, define
\[
 U=C^2-A-jB,\qquad V=\frac{U}{\lambda},\qquad
 b_1=\frac{B/2+jV}{C},\qquad b_0=C-jb_1.
\]
Then the following are equivalent:
\begin{enumerate}
 \item $(b_0+b_1s)^2=A+Bs$;
 \item $C^2V-(B/2+jV)^2=0$;
 \item $\Gamma_j(A,B,C)=0$, where
 \begin{align*}
 \Gamma_j(A,B,C)
 ={}&4\delta C^4
   -4\bigl((j^2+\delta)A+2j\delta B\bigr)C^2\\
   &+\bigl(2jA+(j^2+\delta)B\bigr)^2.
 \end{align*}
\end{enumerate}
The correspondence identifies all nonzero square roots
$b=b_0+b_1s$ satisfying $b_0+jb_1\ne0$ with the nonzero solutions
$C$ of the last equation, by $C=b_0+jb_1$.
\end{proposition}
\begin{proof}
First suppose $b^2=A+Bs$ and $C=b_0+jb_1\ne0$.
Comparing coefficients gives
\[
 A=b_0^2+\delta b_1^2,\qquad B=2b_0b_1,
 \qquad C^2=A+jB+\lambda b_1^2.
\]
Hence $V=b_1^2$ and $Cb_1=B/2+jV$, proving the reconstruction formulas
and the second assertion.

Conversely, suppose the second assertion holds and use the displayed
reconstruction formulas.  Because $C\ne0$, it follows that $b_1^2=V$.
Consequently
\[
 b_0b_1=Cb_1-jb_1^2=B/2
\]
and
\begin{align*}
 b_0^2+\delta b_1^2
 &=C^2-2jCb_1+(j^2+\delta)V\\
 &=C^2-jB-(j^2-\delta)V=A.
\end{align*}
Thus $b^2=A+Bs$, and $b\ne0$ follows from $C=b_0+jb_1\ne0$.
Finally, multiplication of the second assertion by $-4\lambda^2$
gives exactly $\Gamma_j=0$.  This also proves that the two
constructions are inverse to one another on the specified charts.
\end{proof}

For $F=\mathbb Q$ and $\delta=2$, this becomes
\[
 \Gamma_j(A,B,C)=
 8C^4-4\bigl((j^2+2)A+4jB\bigr)C^2
 +\bigl(2jA+(j^2+2)B\bigr)^2.
\]
The root coordinates can equivalently be reconstructed as
\begin{align*}
 b_0&=\frac{-4C^2+2j^2A+j(j^2+2)B}{2(j^2-2)C},\\
 b_1&=\frac{2jC^2-2jA-(j^2+2)B}{2(j^2-2)C}.
\end{align*}
Thus neither $V$, $b_0$, nor $b_1$ consumes an additional variable.

\begin{lemma}[A common chart for finitely many points]
\label{lem:common-quadratic-chart}
Let $b^{(1)},\ldots,b^{(N)}$ be nonzero elements of
$\mathbb Q(\sqrt\delta)$, where $\delta$ is a rational nonsquare.
There exists $j\in\{0,1,\ldots,N\}$ such that
\[
 b^{(r)}_0+j b^{(r)}_1\ne0\qquad(1\leq r\leq N).
\]
\end{lemma}
\begin{proof}
For a fixed nonzero pair $(b_0,b_1)$, the expression $b_0+jb_1$
vanishes for at most one rational $j$.  If $b_1=0$, it never
vanishes.  At most $N$ values are therefore excluded, and the
displayed set has $N+1$ distinct elements.
\end{proof}

\subsection{Application to elliptic doubling}

Let $K=\mathbb Q(\sqrt2)$ and let
\[
 E_D:\ b^2=a^3-Da,\qquad D\in K^\times.
\]
The assumption $D\ne0$ makes this Weierstrass cubic nonsingular.
Write $a=a_0+a_1\sqrt2$ and
$a^3-Da=A+B\sqrt2$ with $A,B\in\mathbb Q$.
An affine point with nonzero ordinate is therefore represented,
on the $j$-chart, by only three rational coordinates
$(a_0,a_1,C)$, subject to $C\ne0$ and $\Gamma_j(A,B,C)=0$.
For fixed $D$, this polynomial has total degree at most six in these
three coordinates: $A,B$ have degree at most three and the leading
terms come from a square of their linear combination.

Reconstruct $b$ as above and put
\[
 h=\frac{3a^2-D}{2b},\qquad
 \xi=h^2-2a,\qquad \eta=h(a-\xi)-b.
\]
These rational functions give the usual doubling formula
$[2](a,b)=(\xi,\eta)$.  Since $b\ne0$, the doubled point is affine.
Conversely, every nonidentity point of $2E_D(K)$ has an affine
preimage with nonzero ordinate: a preimage equal to the identity or
to a point with ordinate zero would double to the identity.
For a list of $N$ such doubled points, choose a halfpoint of each
and then apply Lemma~\ref{lem:common-quadratic-chart}.  Thus the
union of the $N+1$ common charts represents all such lists.

There is no change to the group-theoretic assertions required by the
original construction.  In particular, the identity point must be
handled separately if a surrounding formula permits it.  In the
proposed application, the 211 point occurrences are constrained to be
nonidentity by their original nonvanishing conditions; this condition
is part of the application, not a conclusion of the elimination alone.

\subsection{Guards and finite unions without selector variables}

Every rational function used in a fixed chart can be written as a
quotient of rational polynomials.  Suppose the required denominator
factors in that branch are $d_1,\ldots,d_r$, and set
$\Delta=\prod_{k=1}^r d_k$.  The single equation
\[
 Z\Delta-1=0
\]
with one rational witness $Z$ enforces all the nonvanishing conditions
at once.  After imposing this equation, clearing those denominators
from the remaining branch equations preserves both directions of
the existential assertion.  In particular, the factors $C_i$ must
be guarded; merely imposing $\Gamma_j=0$ at $C_i=0$ would introduce
points outside the displayed chart.

If a branch is given by rational polynomial residuals
$f_{j,1},\ldots,f_{j,r_j}$, including its own equation
$Z\Delta_j-1=0$, set
\[
 H_j=\sum_{k=1}^{r_j}f_{j,k}^2.
\]
Over $\mathbb Q$, the union of the branches is then represented by
\[
 \prod_{j=0}^{N}H_j=0.
\]
All branches may reuse the same list of witness-variable names: a
zero of the product belongs to at least one factor, and that factor
alone must satisfy its denominator guard.  Requiring the guards of
inactive branches would generally destroy completeness and is not
part of this construction.  Additional finite choices, such as the
256 fixed height-class choices, can be included in the same product
over pairs of branch indices without selector variables.

\subsection{Variable count and the scope of the proposed construction}

Using the source inventory of 211 nonidentity point occurrences, the
three-coordinate construction contributes $3\cdot211=633$ rational
witness variables.  With the three-witness semilocal anchor proved in
Proposition~\ref{prop:explicit-semilocal} and the other proposed blocks,
the count is
\[
\begin{array}{lr}
\text{Point coordinates} &633\\
\text{Pole witnesses (42 blocks of five)}&210\\
\text{Anchor witnesses}&3\\
\text{Proxy variables }s,t,z&3\\
\text{The variable }\theta&1\\
\text{The variable }e&1\\
\text{Contact variables }u,v&2\\
\text{Input}&1\\
\text{Common inverse}&1\\\hline
\text{Total}&855.
\end{array}
\]
This is a conditional construction count: the remaining source equations
must all be expressed with precisely these blocks, the advertised
nonvanishing hypotheses must be retained, and the branch polynomials
must be explicitly specified and checked.  The square-root elimination
and common-chart argument above are proved, but they do not alone
certify an entire 855-variable integer-definition polynomial.
The much higher degree resulting from the finite product is an
essential tradeoff of this construction, and must be recorded rather
than hidden behind the variable count.
